\section{Вывод уравнений движения}

\subsection{Координаты маятника}
Координаты груза маятника:
\[
x = u + l\sin\varphi, \qquad y = l\cos\varphi.
\]

\subsection{Кинетическая и потенциальная энергия}
Кинетическая энергия:
\[
T = \frac{m}{2}(\dot{x}^2 + \dot{y}^2) = \frac{m}{2}\left(\dot{u}^2 + 2l\dot{u}\dot{\varphi}\cos\varphi + l^2\dot{\varphi}^2\right).
\]

Потенциальная энергия (нулевой уровень — точка подвеса):
\[
\Pi = mgy = mgl\cos\varphi.
\]

\subsection{Уравнение Лагранжа}
Функция Лагранжа \(L = T - \Pi\). Уравнение Лагранжа для координаты \(\varphi\):
\[
\frac{d}{dt}\left(\frac{\partial L}{\partial \dot{\varphi}}\right) - \frac{\partial L}{\partial \varphi} = 0.
\]

После вычислений получаем:
\[
l\ddot{\varphi} + \ddot{u}\cos\varphi - g\sin\varphi = 0.
\]

\subsection{Линеаризация}
При малых углах \(\varphi \ll 1\): \(\sin\varphi \approx \varphi\), \(\cos\varphi \approx 1\). Линейное уравнение:
\[
\ddot{\varphi} - \frac{g}{l}\varphi = -\frac{\ddot{u}}{l}. \tag{1}
\]

\subsection{Неуправляемый маятник}
При \(\ddot{u} = 0\):
\[
\ddot{\varphi} - \frac{g}{l}\varphi = 0.
\]
Характеристическое уравнение \(\lambda^2 - g/l = 0\) имеет корни \(\lambda_{1,2} = \pm\sqrt{g/l}\). Один корень положительный, следовательно, верхнее положение неустойчиво (седло).